This chapter reviews key literature on mechanical counting and bundling systems, emphasizing intermittent-motion mechanisms, design methods, and practical constraints for low-cost, high-speed and compact size solutions.

\section{Mechanical Indexing and Counting Systems}
Mechanical indexing devices (ratchet-pawl assemblies, Geneva mechanisms, escapements) and cam-driven linkages have long been used for precise, repeatable motion in packaging equipment. Their deterministic timing and robustness make them attractive where electronic control is impractical or costly.

\section{Design Methodologies and Simulation}
Modern machine design combines classical kinematics and dynamics with CAD-based modeling and simulation. Kinematic analysis defines motion profiles and timing; FEA and simple dynamic models inform structural sizing, while rapid prototyping shortens iteration cycles.

\section{Materials and Manufacturing Considerations}
Selection of materials, surface finishes, and tolerances affects wear, friction, and reliability. For small-scale, low-cost machines, a balance is needed between durable components (hardened pins, low-friction bushes) and economical manufacturing methods (sheet metal, simple machining).

\section{Existing Low-Cost Solutions and Gaps}
Existing literature describes semi-automatic incense bundling and counting solutions that use simple fixtures and manual feeds; however, comprehensive studies on fully mechanical, low-cost counting-and-binding machines tailored to small manufacturers are scarce. This gap motivates the present work.

% Related Theory moved to chapter 3